% The rank-one criterion as a point-in-ellipsoid test, drawn at three
% coordinates.  A = diag(4, 9/4, 1) is positive definite, so the locus q = 1 is
% the ellipsoid with semi-axes 2, 3/2, 1, and three vectors b run from the
% origin: (1, 3/4, 0) with q = 1/2, (1, 0, sqrt3/2) with q = 1, (0, 3/2, 1) with
% q = 2.  Each of the three lies in a coordinate plane, so the point where its
% ray meets the ellipsoid -- the open circle, at 1/sqrt(q) times b -- sits on a
% drawn principal ellipse.  Orthographic projection of R^3 with azimuth 20 and
% elevation 22 degrees, scaled by 2.20cm:
%   X = 2.20(-0.342020 x_1 + 0.939693 x_2)
%   Y = 2.20(-0.352015 x_1 - 0.128123 x_2 + 0.927184 x_3)
% Every coordinate below is the parenthesised value, with its source triple in a
% trailing comment; the 2.20cm is carried by the x and y unit vectors.  A curve
% is given by the projections of two conjugate semi-axis vectors, so that
% (X, Y) = (cos u) proj(first) + (sin u) proj(second).
\begin{tikzpicture}[
  x=2.20cm, y=2.20cm,
  vec/.style={-{Latex[length=2mm,width=1.5mm]}, line width=0.75pt},
  frame/.style={densely dotted, line width=0.3pt, black!35},
  edge/.style={line width=0.45pt, black!65},
  rim/.style={line width=0.45pt, black!70},
  hidden/.style={line width=0.45pt, black!45, dashed, dash pattern=on 1.4pt off 1.2pt},
  slack/.style={densely dotted, line width=0.5pt, black!55},
  hit/.style={line width=0.5pt, black!85},
  lead/.style={line width=0.3pt, black!55},
  every node/.style={font=\small}]

% ---- the three coordinate axes, as a frame of reference ------------------
% each runs from the origin past its pole, which sits on the ellipsoid
\draw[frame] (0,0) -- (-0.807167,-0.830755);   % (2.36, 0,     0    )
\draw[frame] (0,0) -- ( 1.620970,-0.221012);   % (0,    1.725, 0    )
\draw[frame] (0,0) -- ( 0,        1.038446);   % (0,    0,     1.12 )

% ---- the ellipsoid: three principal ellipses, hidden arcs first ----------
% A point x of the ellipsoid faces the viewer when (A^{-1}d).x > 0, with
% d = (0.871268, 0.317116, 0.374607) the viewing direction; that inequality
% cuts each principal ellipse at the two parameters recorded below.
% E_12, semi-axis vectors (2,0,0) -> (-0.684040,-0.704030)
%                     and (0,3/2,0) -> ( 1.409539,-0.192185)
\draw[hidden, variable=\u, domain=115.887:295.887, samples=61, smooth]
  plot ({-0.684040*cos(\u)+1.409539*sin(\u)}, {-0.704030*cos(\u)-0.192185*sin(\u)});
% E_13, semi-axis vectors (2,0,0) and (0,0,1) -> (0, 0.927184)
\draw[hidden, variable=\u, domain=130.693:310.693, samples=61, smooth]
  plot ({-0.684040*cos(\u)}, {-0.704030*cos(\u)+0.927184*sin(\u)});
% E_23, semi-axis vectors (0,3/2,0) and (0,0,1)
\draw[hidden, variable=\u, domain=150.562:330.562, samples=61, smooth]
  plot ({1.409539*cos(\u)}, {-0.192185*cos(\u)+0.927184*sin(\u)});
\draw[edge, variable=\u, domain=-64.113:115.887, samples=61, smooth]
  plot ({-0.684040*cos(\u)+1.409539*sin(\u)}, {-0.704030*cos(\u)-0.192185*sin(\u)});
\draw[edge, variable=\u, domain=-49.307:130.693, samples=61, smooth]
  plot ({-0.684040*cos(\u)}, {-0.704030*cos(\u)+0.927184*sin(\u)});
\draw[edge, variable=\u, domain=-29.438:150.562, samples=61, smooth]
  plot ({1.409539*cos(\u)}, {-0.192185*cos(\u)+0.927184*sin(\u)});
% the silhouette: the contour generator is the section of the ellipsoid by the
% plane (A^{-1}d).x = 0, and its projection is the ellipse through
% (1.566752, 0.134478) and (0, 1.172254) at the two conjugate parameters,
% equivalently the ellipse of semi-axes 1.579547 and 1.162758 turned by
% 10.8171 degrees
\draw[rim, rotate=10.8171] (0,0) ellipse (1.579547 and 1.162758);

% ---- the three vectors, each with the point where its ray leaves ----------
% q = 1/2: the ray leaves at sqrt2 times b, on E_12 at parameter 45 degrees
\draw[vec]   (0,0) -- (0.362749,-0.448107);                 % b = (1, 3/4, 0)
\draw[slack] (0.362749,-0.448107) -- (0.513005,-0.633719);  % (1.414214, 1.060660, 0)
\draw[hit]   (0.513005,-0.633719) circle (1.7pt);
% q = 1: the tip is itself the exit point, on E_13 at parameter 60 degrees
\draw[vec] (0,0) -- (-0.342020,0.450950);                   % b = (1, 0, sqrt3/2)
\draw[hit] (-0.342020,0.450950) circle (1.7pt);
% q = 2: the shaft leaves at b over sqrt2, on E_23 at parameter 45 degrees
\draw[vec] (0,0) -- (1.409539,0.734999);                    % b = (0, 3/2, 1)
\draw[hit] (0.996695,0.519723) circle (1.7pt);              % (0, 1.060660, 0.707107)
\fill (0,0) circle (1.1pt);

% ---- labels ---------------------------------------------------------------
\node[anchor=north east, font=\footnotesize] at (-0.83,-0.85) {$b_1$};
\node[anchor=west,       font=\footnotesize] at ( 1.66,-0.23) {$b_2$};
\node[anchor=east,       font=\footnotesize] at (-0.04, 1.05) {$b_3$};
\draw[lead] (-0.60,1.140) -- (-0.536,1.056);   % to the rim at parameter 110 degrees
\node[anchor=east, align=right, font=\footnotesize] at (-0.62,1.19)
  {$b\adj A^{-1}b=1$};
\node[align=center, font=\footnotesize] at ( 0.80,-0.31) {$q=\tfrac12$\\[-1pt] inside};
\node[align=center, font=\footnotesize] at (-0.74, 0.36) {$q=1$\\[-1pt] on};
\node[anchor=west, align=left, font=\footnotesize] at ( 1.46, 0.80)
  {$q=2$\\[-1pt] outside};

% ---- the reading, below the picture ---------------------------------------
\node[anchor=north, align=left, font=\scriptsize] at (0.14,-1.50) {%
  $A=\diag(4,\tfrac94,1)$,\quad $q=b\adj A^{-1}b$,\quad
  semi-axes $\sqrt{A_{11}},\sqrt{A_{22}},\sqrt{A_{33}}=2,\tfrac32,1$\\[4pt]
  \begin{tabular}{@{}l@{\hspace{1.3em}}l@{\hspace{1.3em}}l@{\hspace{1.3em}}l@{}}
    $b=(1,\tfrac34,0)$          & $q=\tfrac12$ & $A-bb\adj\pd0$
      & $C$ dominates strictly\\[1pt]
    $b=(1,0,\tfrac{\sqrt3}{2})$ & $q=1$        & $A-bb\adj\psd0$, singular
      & $C$ dominates, not strictly\\[1pt]
    $b=(0,\tfrac32,1)$          & $q=2$        & $A-bb\adj$ indefinite
      & $C$ does not dominate\\
  \end{tabular}\\[4pt]
  at $m=k+1$, with $C=\{1,\dots,m\}\setminus\{c_0\}$: \
  $A=\diag(1-t_c)_{c\in C}$, \ $b=(z_c)_{c\in C}$, \ $q=\sum_{c\ne c_0}p_c$};

\end{tikzpicture}
